% ===========================================================================
\section{Divergence: comparing two laws}\label{sec:divergence}
\warmth{0}\quad\whisper{KL divergence is not a distance, but distance intuition is useful.}

\begin{strand}
$\KL(p\Vert q)$ measures the inefficiency that arises when the true law is $p$
but we assume the law is $q$. It is asymmetric, so it is not a distance between
distributions in the strict mathematical sense. Nevertheless, treating it as a
directed notion of ``distance'' is often very useful.
\end{strand}

\begin{definition}[Kullback--Leibler divergence]\label{def:kl}
For pmfs $p$ and $q$ on $\X$,
\[
  \KL(p\Vert q)=\sum_{x\in\X}p(x)\log\frac{p(x)}{q(x)}.
\]
We take $0\log(0/q)=0$. If $p(x)>0$ for some $x$ with $q(x)=0$, then
$\KL(p\Vert q)=\infty$.
\end{definition}

\block{Expectation and entropy forms}
If $X\sim p$, then
\begin{align*}
  \KL(p\Vert q)
  &=\E_{X\sim p}\!\left[\log\frac{p(X)}{q(X)}\right]\\
  &=-\E_{X\sim p}[\log q(X)]-\Ent(X).
\end{align*}

\begin{yourturn}
Fill the algebraic step:
\[
  \E_p\!\left[\log\frac{p(X)}{q(X)}\right]
  =\E_p[\log p(X)]-\E_p[\log q(X)]
  =\answerin[4.6cm]{-\Ent(X)-\E_p[\log q(X)]}.
\]
\recall{$-\Ent(X)-\E_p[\log q(X)]$.}
\end{yourturn}

\block{The order cannot be reversed}

\Needspace{11\baselineskip}
\begin{example}[asymmetry]
On $\X=\{0,1\}$, let
\[
  p=(1/2,1/2),\qquad q=(1,0).
\]
Then $\KL(p\Vert q)=\infty$ because $q(1)=0$ while $p(1)>0$, but
\[
  \KL(q\Vert p)=\log\frac{1}{1/2}=1\ \text{bit}.
\]
Thus $\KL(p\Vert q)$ and $\KL(q\Vert p)$ need not agree.
\end{example}

\begin{yourturn}
Explain in one sentence why the first direction is infinite but the reverse
direction is finite.
\workspace[2]
\answerprose{$\KL(p\Vert q)$ sums $p(x)\log\frac{p(x)}{q(x)}$; since $p(1)>0$ while $q(1)=0$, the term $\log\frac{p(1)}{q(1)}=\log\infty$ blows up. In the reverse direction every $x$ with $q(x)>0$ also has $p(x)>0$, so no infinite term appears.}
\recall{$p(1)>0=q(1)$ makes $\KL(p\Vert q)=\infty$; the reversed support has no such term.}
\end{yourturn}

\block{Comparison with the uniform law}

\begin{proposition}[uniform reference]\label{prop:uniform-kl}
If $q$ is uniform on a finite alphabet $\X$, then for $X\sim p$,
\[
  \KL(p\Vert q)=\log|\X|-\Ent(X).
\]
\end{proposition}
\begin{proof}
Since $q(x)=1/|\X|$,
\begin{align*}
  \KL(p\Vert q)
  &=\sum_xp(x)\log p(x)+\sum_xp(x)\log|\X|\\
  &=-\Ent(X)+\log|\X|\sum_xp(x)\\
  &=\log|\X|-\Ent(X),
\end{align*}
where the last equality uses $\sum_xp(x)=1$.
\end{proof}

\begin{yourturn}
Recall the key simplification for a uniform reference law:
\[
  \KL(p\Vert q)
  =-\Ent(X)+\log|\X|\sum_xp(x)
  =\answerin[3.2cm]{\log|\X|-\Ent(X)}.
\]
\recall{$\log|\X|-\Ent(X)$.}
\end{yourturn}
